\documentclass[tikz,multi]{standalone}
%\usepackage
%\usepackage{tikz}
\usepackage{luacode}

\begin{luacode*}
N=8
SCALE = 12

A={0,0}
B={1,0}
C = {1/2,math.sqrt(3)/2}

function lerp(P,Q,t)
	return {P[1]*(1-t) + Q[1] * t, P[2]*(1-t) + Q[2] * t}
end



function sierpinski(P,Q,R,n)
	if n == 0 then
		printTriangle(P,Q,R)
	else
		sierpinski(P, lerp(P,Q,1/2), lerp(P,R,1/2), n-1)
		sierpinski(lerp(Q,P,1/2), Q, lerp(Q,R,1/2), n-1)
		sierpinski(lerp(R,P,1/2), lerp(R,Q,1/2), R, n-1)
	end
end

function printTriangleWithTikz(P,Q,R)
tex.print(string.format(
			"\\fill[%s] (%f,%f) -- (%f,%f) -- (%f,%f) -- cycle;",
			"black",
			P[1], P[2],
			Q[1], Q[2],
			R[1], R[2]
		))
end

printTriangle=printTriangleWithTikz

POINTS_PER_CM = 28.35
COEF = SCALE * POINTS_PER_CM

function printTriangleWithPdfliteral(P,Q,R)
tex.print(string.format(
			"\\leavevmode\\hbox{\\pdfextension literal direct {q 0 0 0 rg %f %f m %f %f l %f %f l h f Q}}",
			COEF * P[1], COEF * P[2],
			COEF * Q[1], COEF * Q[2],
			COEF * R[1], COEF * R[2]
		))
end

t0 = os.clock()
\end{luacode*}

\begin{document}
\begin{tikzpicture}[scale=\directlua{tex.print(SCALE)}]
\directlua{sierpinski(A, B, C,N)}
\end{tikzpicture}

\newpage

\begin{luacode*}
t1=os.clock()
printTriangle=printTriangleWithPdfliteral
sierpinski(A, B, C, N)
\end{luacode*}

\newpage

Statistics:

Depth : \directlua{tex.print(N)}

Number of triangles :  $3^N=\directlua{tex.print(math.tointeger(3^N))}$.


TikZ picture printed in \directlua{tex.print(t1-t0)} seconds.



Pdfliteral picture printed in \directlua{t2=os.clock();tex.print(t2-t1)} seconds.





\end{document}
